%قبل از شروع مسئله ترجیحاً چیزی ننویسید. اما اگر خواستید بنویسید، در قالب لیستی (فقرات و شمارش) بنویسید.
%برای اضافه کردن نمره، بعدی از شروع‌مسئله، [x] را اضافه کنید و به جای x  نمره مورد نظر را قرار دهید. برای زیرمسئله‌ها هم به همین صورت است.
%\شروع‌مسئله[10]
\شروع‌مسئله[20] \متن‌سیاه{تحلیل مجانبی}. در هر قسمت علاوه بر پاسخ درست، روابطی که استفاده کردید را نوشته و به طور مختصر توضیح دهید.
% صورت اصلی مسئله را در این‌جا بنویسید.
\شروع{زیرمسائل}
\زیرمسئله[5] توابع زیر را از لحاظ مجانبی مرتب کنید.
$$f_1(n) = (2n)!, f_2(n) = \binom{2n}{n}, f_3(n) = 2^{2n}, f_4(n) = (2n)^n, f_5(n)=n^{2n}$$

\پاسخ{
	ترتیب توابع از کوچک به بزرگ و از چپ به راست:
	$$(f_2, f_3, f_4, f_1, f_5)$$
	این ترتیب با استفاده از تقریب استرلینگ به راحتی به دست می‌آید. به طور مثال باید $f_2$ را به صورت $\frac{(2n)!}{n!\cdot n!}$ نوشت و سپس هر کدام از عبارت‌های فاکتوریل را با تقریب استرلینگ ساده کرد. با این کار می‌توان نشان داد $f_2$ از بقیه توابع (به صورت مجانبی) کوچک‌تر است. همچنین $f_1$ بین $f_4$ و $f_5$ قرار دارد. مشابه این سؤال در تمرین اول داده شده است.
	
}

\نمره‌دهی{
	اگر ترتیب درست باشد کل نمره داده می‌شود. اگر ترتیب ۲ تابع اشتباه باشد یک نمره کسر می‌شود. برای ترتیب اشتباه ۳ تابع، ۳ نمره و برای ترتیب اشتباه ۴ تابع ۴ نمره کم می‌شود. ترتیب اشتباه ۵ تابع هم نمره‌ای نمی‌گیرد.
	
}

\زیرمسئله[5] رابطۀ بازگشتی زیر را حل کنید. (راهنمایی: از تغییر متغیر $n=2^m$ استفاده کنید)
\begin{center}
$T(n)=4T(\sqrt{n})+\log^2 n$
\end{center}

\پاسخ{
	در کلاس گفتیم اگر مبنای $\log$ گفته نشده، در مبنای ۲ هست. بنابراین با تغییر متغیر گفته شده و تعریف تابع $S$ داریم.
	$$T(2^m) = 4T(2^{m/2}) + m^2 \Rightarrow S(m) = 4S(m/2) + m^2$$
	بنابر قضیۀ اصلی داریم
	$$S(m) = \Theta(m^2\log m) \Rightarrow  T(2^m) = \Theta(m^2\log m) = \Theta(\log^2 n \cdot \log\log n).$$
	
	مشابه این سؤال در تمرین دوم داده شده است.
	
}

\نمره‌دهی{
	اگر جواب نهایی به صورت $\log^2 n \log\log n$ یا $m^2\log m$ نوشته شود نمره کامل و در غیر این صورت با توجه به راه‌حل نمره داده می‌شود.
	
}
%%2011FQ1

%\زیرمسئله[4] رابطۀ بازگشتی زیر را حل کنید.
%$$T(n)=\log n + T(\sqrt{n})$$
%%2011FQ1

\زیرمسئله[5]
با استفاده از قضیۀ اصلی، رابطۀ زیر را حل کنید.
\begin{center}
$T(n) = 4T(n/2)+\Theta(n^2 \log^2 n)$
\end{center}
%sol n^2 log^3 n
\پاسخ{
	با توجه به قضیۀ اصلی (صورتی که در جلسۀ سوم کلاس مطرح شد) حالت دوم (تساوی) برقرار است (به طور شهودی می‌توان گفت عبارت‌های لگاریتمی تساوی را به هم نمی‌زنند). بنابر این $T(n) = \Theta(n^2\log^3 n)$ است. حل روابط به قضیۀ اصلی در تمرین دوم مطرح شده است.
	
}
\نمره‌دهی{
	اگر جواب نهایی درست باشد نمره کامل وگرنه با توجه به راه‌حل نمره داده خواهد شد: یک نمره برای به‌دست آوردن $L$ و ۴ نمره برای مقایسه و به‌دست آوردن $\Theta$.
	
}


\زیرمسئله[5] برای مسئله‌ای که طول ورودی آن $n$ و طول خروجی آن $d$ ($1\leq d\leq n$) است دو الگوریتم داریم. زمان اجرای الگوریتم اول $\Theta(n+d^2)$ و زمان اجرای الگوریتم دوم 
$\Theta(\delta\cdot n^{2-\frac 2 7})$
 است که $\delta=\frac n d$. الگوریتم سومی به این صورت طراحی می‌کنیم که الگوریتم اول و دوم را همزمان اجرا می‌کند و هر کدام زودتر به پایان رسید، جواب آن را برمی‌گرداند. زمان اجرای الگوریتم سوم را در بدترین حالت، فقط برحسب $n$ بیان کنید. (راهنمایی: حالتی را درنظر بگیرید که زمان اجرای دو الگوریتم با هم برابر می‌شود)

\پاسخ{
وقتی $d=1$ است الگوریتم اول بهتر است و وقتی $d=n$ باشد الگوریتم دوم بهتر است. همچنین در الگوریتم اول زمان با افزایش $d$ افزایش و در الگوریتم دوم زمان با افزایش $d$ کاهش می‌یابد. بنابراین بدترین حالت وقتی است که زمان اجرای دو الگوریتم با هم برابر شود. در این حالت با صرف‌نظر از $n$ در زمان اجرای الگوریتم اول داریم
$$d^2 = n^{3-(2/7)} / d \Rightarrow d^3 = n^{3-(2/7)} \Rightarrow d = n^{1-(2/21)}.$$
در این حالت زمان اجرای الگوریتم برابر است با $O(n^{2-4/21})$ و برای بقیه $d$ها از این مقدار کمتر است.
از این رابطه در یکی از لم‌های مقاله زیر استفاده شده است:
\begin{latin}
Boroujeni, M., Ehsani, S., Ghodsi, M., HajiAghayi, M. and Seddighin, S., 2018. Approximating edit distance in truly subquadratic time: Quantum and MapReduce. In SODA (pp. 1170-1189). SIAM.
\end{latin}

}
\نمره‌دهی{
با توجه به راه‌حل و جواب نهایی نمره داده می‌شود. راه‌حل درست ۴ نمره و جواب درست ۱ نمره دارد.

}

%\زیرمسئله[4]
%الگوریتمی داریم که زمان اجرای آن
%$O(n^{2-(3/2)\phi} + n^{2\beta} + n^{(1-\beta)(2-\phi)+1.5\beta})$
%است. پارامترهای $\beta$ و $\phi$ را می‌توانیم به دلخواه از اعداد بین صفر تا یک انتخاب کنیم. هدف این که با انتخاب پارامترهای مناسب زمان اجرای الگوریتم را تا حد ممکن کاهش دهیم. پارامترها مناسب و زمان اجرای الگوریتم با این پارامترها چقدر خواهد بود؟ (راهنمایی: زمانی الگوریتم کمترین زمان را دارد که توان سه جمله یکسان باشد)



%اگر مسئله چند قسمت داشت به این صورت عمل کنید:
%\شروع{زیرمسائل}
%\زیرمسئله[0] % مسئلۀ ۱آ
%\زیرمسئله % مسئلۀ ۱ب
%\زیرمسئله % مسئلۀ ۱پ
%\زیرمسئله % مسئلۀ ۱ت
\پایان{زیرمسائل}




%%%%%%%%%%%%%%%%%%%%%%%%%%%%%%%%%%%%%%%%%%%%%%%%%%%%%
% برخی دستورات مفید برای تمرین‌ها. %
%%%%%%%%%%%%%%%%%%%%%%%%%%%%%%%%%%%%%%%%%%%%%%%%%%%%%

% فرمول‌های نمونه:
%$f(x) = \Theta(x)$
%$T(x, y) \leq \log(x) + 2^y + \binom{2n}{n}$
% {\tt code\_function}


% لیست موارد بدون شماره‌گذاری:
%\شروع{فقرات}
%    \فقره مورد اول 
%        \شروع{فقرات}
%            \فقره مورد اول داخلی 
%                \شروع{فقرات}
%                    \فقره اولین مورد خیلی داخلی 
%                    \فقره دومین مورد خیلی داخلی 
%                \پایان{فقرات}
%            \فقره مورد دوم داخلی 
%        \پایان{فقرات}
%    \فقره مورد دوم 
%\پایان{فقرات}

% لیست موارد شماره‌گذاری شده:
%\شروع{شمارش}
%    \فقره یک مورد 
%    \فقره یک مورد دیگر 
%    \فقره باز هم یک مورد دیگر
%\پایان{شمارش}

% برای نوشتن یک فرمول چند خطی که یک شماره بخوره. برای حالت‌های دیگر به سایت زیر مراجعه کنید:
% https://www.overleaf.com/learn/latex/Aligning_equations_with_amsmath
%\begin{align} 
%    \begin{split}
%        (x+y)^3 &= (x+y)^2(x+y) \\
%                &= (x^2+2xy+y^2)(x+y) \\
%                &= (x^3+2x^2y+xy^2) + (x^2y+2xy^2+y^3) \\
%                &= x^3+3x^2y+3xy^2+y^3
%    \end{split}                                 
%\end{align}

% نوشتن ماتریس و آرایه دوبعدی
%\begin{align}
%    A = 
%    \begin{bmatrix}
%        A_{11} & A_{21} \\
%        A_{21} & A_{22}
%    \end{bmatrix}
%\end{align}

% اضافه کردن تصویر با شرح. برای اضافه کردن تصویر بدون شرح، فقط دستور درج‌تصویر را استفاده کنید.
%\شروع{شکل}[h!]
%\تنظیم‌ازوسط
%  \درج‌تصویر[پهنا=0.5\پهنای‌متن]{img.jpg}
%\شرح{شرح تصویر}
%\برچسب{تصویر:۱}
%\پایان{شکل}
%\پایان{زیرمسائل}

%کد زیر را ببینید.
%%\lstinputlisting{source_filename.py}
%\begin{latin}
%\begin{lstlisting}
%# Import the modules
%import sys
%import random
%
%ans = True
%
%while ans:
%question = raw_input("Ask the magic 8 ball a question: (press enter to quit) ")
%
%answers = random.randint(1,8)
%
%if question == "":
%sys.exit()
%
%elif answers == 1:
%print "It is certain"
%
%elif answers == 2:
%print "Outlook good"
%
%elif answers == 3:
%print "You may rely on it"
%
%elif answers == 4:
%print "Ask again later"
%
%elif answers == 5:
%print "Concentrate and ask again"
%
%elif answers == 6:
%print "Reply hazy, try again"
%
%elif answers == 7:
%print "My reply is no"
%
%elif answers == 8:
%print "My sources say no"
%\end{lstlisting}
%\end{latin}
